\section{Conclusão}

\paragraph{}
As atividades desenvolvidas permitiram a aplicação prática de técnicas de
\textit{Open Source Intelligence} (OSINT), demonstrando como informações
disponíveis publicamente podem ser coletadas, correlacionadas e analisadas para
apoiar processos de investigação e segurança da informação.

\paragraph{}
Durante os exercícios foram utilizadas diversas ferramentas e fontes abertas,
incluindo consultas DNS com \texttt{dig}, pesquisas em bancos de dados públicos
por meio do Shodan, análises de vazamentos de dados utilizando o Maltego e o
plugin \textit{Have I Been Pwned?}, além da utilização do Google Hacking
Database (GHDB) e de Google Dorks para identificação de informações
potencialmente expostas na Internet.

\paragraph{}
Os resultados obtidos evidenciam que informações aparentemente inofensivas,
quando analisadas em conjunto, podem revelar detalhes relevantes sobre
infraestruturas, serviços, diretórios expostos, vazamentos de dados e outros
ativos digitais. Esse processo de correlação de informações é uma etapa
fundamental em atividades de reconhecimento e pode auxiliar tanto equipes de
segurança defensiva quanto profissionais envolvidos em testes de penetração.

\paragraph{}
A realização das atividades também permitiu compreender a importância da
proteção de informações públicas, da configuração adequada de serviços
acessíveis pela Internet e do monitoramento contínuo de vazamentos de dados,
reduzindo riscos associados à exposição indevida de informações sensíveis.

\paragraph{}
Por fim, a prática reforçou o papel do OSINT como uma das fases mais
importantes do ciclo de segurança ofensiva, fornecendo subsídios para análises
mais aprofundadas e contribuindo para uma melhor compreensão da superfície de
ataque de organizações e sistemas.
